\documentclass[../main]{subfiles}
\begin{document}

\chapter{Máquinas Sincrónicas}

\section{Introducción a las Máquinas Sincrónicas}

Las máquinas sincrónicas son dispositivos electromecánicos que convierten energía eléctrica en energía mecánica y viceversa. Este tipo de máquinas se utilizan principalmente en aplicaciones donde se requiere un control preciso de la velocidad y el torque\footnote{Momento de fuerza} como en generadores y motores industriales. 

\info{Principio de Funcionamiento}{
El principio básico de funcionamiento de una máquina sincrónica se basa en la interacción entre un campo magnético giratorio, creado por el estator\footnote{Parte estacionaria de una máquina eléctrica} y el rotor\footnote{Parte giratoria de una máquina eléctrica}. El rotor gira a la misma velocidad que el campo magnético del estator, de ahí el término "sincrónico".

\begin{equation}
n_s = \frac{120 \cdot f}{P}
\end{equation}

Donde:

\begin{itemize}
    \item $n_s$: Velocidad sincrónica (RPM)
    \item $f$: Frecuencia de la corriente (Hz)
    \item $P$: Número de polos
\end{itemize}

Esta fórmula es crucial para determinar la velocidad de rotación de la máquina sincrónica en función de la frecuencia de la corriente y el número de polos de la máquina.
}

\info{Impedancia Síncrona}{
La impedancia sincrónica es una medida de la oposición total que una máquina sincrónica presenta al flujo de corriente alterna. Se expresa en términos de resistencia y reactancia\footnote{Oposición que presentan las bobinas y condensadores al paso de corriente alterna}. La relación de la impedancia sincrónica se puede expresar como:

\begin{equation}
Z_s = R_a + jX_s
\end{equation}

Donde:

\begin{itemize}
    \item $Z_s$: Impedancia sincrónica
    \item $R_a$: Resistencia del devanado del estator
    \item $X_s$: Reactancia sincrónica
    \item $j$: Unidad imaginaria
\end{itemize}

Esta fórmula es útil para analizar el comportamiento de la máquina bajo diferentes condiciones de carga.
}

\info{Regulación de Voltaje}{
La regulación de voltaje en una máquina sincrónica es el cambio en el voltaje terminal cuando la máquina pasa de una condición de carga plena a una condición sin carga, manteniendo constante la velocidad y la excitación. Se puede calcular mediante la fórmula:

\begin{equation}
\text{Regulación de Voltaje (\%)} = \frac{V_{NL} - V_{FL}}{V_{FL}} \times 100
\end{equation}

Donde:

\begin{itemize}
    \item $V_{NL}$: Voltaje sin carga
    \item $V_{FL}$: Voltaje a plena carga
\end{itemize}
}

\ejemplo{Ejemplo}{Calcular la velocidad sincrónica de una máquina sincrónica que tiene 6 polos y opera a una frecuencia de 50 Hz.

Datos:

$f = 50 \text{ Hz}$

$P = 6$

\begin{equation}
n_s = \frac{120 \cdot 50}{6} = 1000 \text{ RPM}
\end{equation}

Por lo tanto, la velocidad sincrónica de la máquina es 1000 RPM.
}

\actividad{Responder el siguiente cuestionario}
\begin{enumerate}
    \item ¿Qué es una máquina sincrónica y cuáles son sus principales aplicaciones?
    \item ¿Cómo se calcula la velocidad sincrónica de una máquina sincrónica?
    \item ¿Qué es la impedancia sincrónica y cuáles son sus componentes?
    \item ¿Qué significa la regulación de voltaje en una máquina sincrónica?
    \item ¿Cuáles son las diferencias entre un motor sincrónico y un motor de inducción?
    \item ¿Qué factores afectan la estabilidad de una máquina sincrónica?
    \item ¿Cómo se logra el arranque de una máquina sincrónica?
    \item ¿Qué es la reactancia sincrónica y cómo afecta el rendimiento de la máquina?
    \item ¿Qué es el factor de potencia y cómo se relaciona con las máquinas sincrónicas?
    \item ¿Cuáles son las ventajas y desventajas de utilizar máquinas sincrónicas en la industria?
\end{enumerate}

\actividad{Resolver los siguientes problemas}
\begin{enumerate}
    \item Calcular la velocidad sincrónica de una máquina que tiene 4 polos y opera a una frecuencia de 60 Hz.
    \item Determinar la impedancia sincrónica de una máquina si la resistencia del devanado es 0.5 ohmios y la reactancia sincrónica es 3 ohmios.
    \item Calcular la regulación de voltaje de una máquina sincrónica si el voltaje sin carga es 240 V y el voltaje a plena carga es 220 V.
    \item Hallar la velocidad sincrónica de una máquina que opera a 75 Hz y tiene 8 polos.
    \item Calcular la impedancia sincrónica total si $R_a = 1 \Omega$ y $X_s = 4 \Omega$.
    \item Determinar la velocidad sincrónica para una máquina de 10 polos que opera a una frecuencia de 50 Hz.
    \item Calcular la regulación de voltaje si $V_{NL} = 230 V$ y $V_{FL} = 210 V$.
    \item Hallar la velocidad sincrónica de una máquina sincrónica con 2 polos operando a 60 Hz.
    \item Calcular la impedancia sincrónica si la resistencia es 0.8 ohmios y la reactancia es 2.5 ohmios.
    \item Determinar la velocidad sincrónica para una máquina de 12 polos que opera a una frecuencia de 60 Hz.
\end{enumerate}

\end{document}